\chapter{Monte Carlo Evidence}
\label{ch:mc}

On synthetic panels where the truth, the $\theta$ ordering, is known, both covariance
constructions attain nominal coverage, and the direct construction remains calibrated
where the MDS reconstruction can fail. The two constructions of Chapter~\ref{ch:cross}
coincide on a balanced panel, exactly under a common bandwidth and asymptotically under
per-series bandwidths (Proposition~\ref{prop:equivalence}), but target different objects
on an unbalanced one (Remark~\ref{rem:unbalanced-targets}). The table script writes
Tables~\ref{tab:mc-coverage} and~\ref{tab:mc-crossover} and the tables of
Appendix~\ref{app:add-mc} in a single run, using the designs of the notebook
\texttt{01\_toydataset}. The figure script draws Figure~\ref{fig:mc-crossover}. The NGDP
draw diagnostics of \S\ref{sec:mc-mechanism} come from the last section of the notebook
(Table~\ref{tab:impl-repro}).

\section{Design}
\label{sec:mc-design}

The data-generating process is the additive common-shock model the theory assumes,
\begin{equation}
  X_{t,j} \;=\; \theta_j + c_t + w_t\,u_{t,j},
  \qquad t=1,\dots,T,\ \ j=1,\dots,p,
  \label{eq:dgp}
\end{equation}
with four controllable ingredients. The idiosyncratic term $u_{t,j}$ is AR(1) with
coefficient $\rho$ (serial dependence). The common shock $c_t$ is additive with variance
$\sigma_c^2$ (cross-sectional dependence) and cancels in the differences
$d^{jk}_t = \theta_j-\theta_k + w_t(u_{t,j}-u_{t,k})$. The volatility path $w_t$ has a
few localized episodes (non-stationarity, crisis quarters). Staggered entry and exit
introduce panel imbalance. The means $\theta_j$ fix the true ranks (rank~1 is the
smallest $\theta$). Under an imbalance $\iota$, forecaster $j$ is observed on a single
window of $\operatorname{round}((1-\iota)T)$ consecutive dates, with the windows' starts
spread evenly over the sample. The non-stationary designs set $w_t$ to the episode height
(\texttt{vol\_scale}, $12$ unless stated) on five windows of eight dates at random
positions and $w_t=1$ elsewhere; only the idiosyncratic part is scaled.

For each design I draw $R=500$ replications and run the stepwise simulation procedure at
$\alpha=0.05$ with $B=2000$ draws under both covariance constructions on the same panels
(a paired comparison). I record the empirical joint coverage
of the rank sets ($\Prob$ that every true rank lies in its confidence set), the mean
rank-set width (lower is sharper), the false-separation rate on a design with tied
$\theta$, and the projection distance $\Frob{\Sh-\Sh^{+}}$ or
$\Frob{\Omegahat-\Omegahat^{+}}$. The block bootstrap of \S\ref{par:mbb} runs on the
same panels, with $B=2000$ resamples, as a covariance-free benchmark.

\section{Coverage and equivalence}
\label{sec:mc-coverage}

Table~\ref{tab:mc-coverage} compares the two constructions across sweeps around a
baseline ($p=5$, $T=200$, $\rho=0.3$, $\sigma_c=1$, balanced) and on the combined
realistic design. Three facts stand out. First, both constructions hold nominal joint
coverage everywhere (never below $0.968$ at the nominal $0.95$), so the procedure is
valid with either covariance input. Second, on the balanced stationary designs the two
are close to equivalent. The MDS projection distance is numerically zero and the widths
agree within $0.01$, except at $\rho=0.9$, where MDS is $0.03$ sharper.
Proposition~\ref{prop:equivalence}(iii) predicts exactly this: under the per-series
Andrews bandwidths used here, the two constructions agree only asymptotically. The same
fact explains the small positive projection distance of $\Omegahat$ on balanced panels.
Third, the common shock is inert. It cancels exactly in every $d^{jk}_t$, so the
$\sigma_c$ designs differ from the baseline only through their seeds, and the spread of
these rows ($0.44$--$0.47$) indicates the size of the Monte Carlo noise in the widths.

The block bootstrap covers at or above nominal on every
design (never below $0.966$). On the balanced stationary designs its widths are within
$0.02$ of $\OmegaDirect$'s, except at $\rho=0.9$, where it is the sharpest of the three
($1.96$ against $2.00$ for MDS and $2.02$ for $\OmegaDirect$) and covers about as often
($0.966$ against $0.968$; Appendix~\ref{app:add-mc}). Under imbalance alone it is the
widest of the three ($1.107$ against $1.050$ for $\OmegaDirect$ and $1.005$ for MDS at
imbalance $0.4$). On the realistic design it is sharper than MDS ($3.566$ vs.\ $3.765$)
and sides with $\OmegaDirect$ ($3.518$). Because the bootstrap uses no covariance
estimate, this pattern is consistent with the MDS widths there reflecting lost
calibration rather than the data.

\tabMcCoverage

On the tie design (a cluster of equal-$\theta$ forecasters), the false-separation rate
per tied pair is $0.02$ for both constructions and for the block bootstrap, below the
nominal $\alpha=0.05$, so no route confidently orders forecasters that are
genuinely equal. Appendix~\ref{app:add-mc} reports the full grid of designs, with
marginal coverage.

\section{The calibration mechanism, and where the direct construction wins}
\label{sec:mc-mechanism}

Imbalance alone (row ``imbalance 0.4'') keeps MDS marginally sharper ($1.005$ vs.\
$1.050$), and non-stationarity alone (balanced) is nearly neutral. Their combination, the
realistic design, raises the MDS projection distance from essentially zero to $0.25$ and
flips the width ordering: $\OmegaDirect$ becomes sharper ($3.518$ vs.\ $3.765$). Under
imbalance the MDS route also targets a different object from $\OmegaDirect$
(Remark~\ref{rem:unbalanced-targets}), but on the stationary imbalance designs that
misspecification costs little. The cost appears when it is compounded by a loss of
calibration in the MDS route.

Table~\ref{tab:mc-crossover} makes the crossover explicit. As the height of the
volatility episodes (\texttt{vol\_scale}) increases with imbalance held fixed, the MDS
projection distance grows monotonically, and the MDS sets become wider than
$\OmegaDirect$'s from a height of $6$ onward. Coverage stays at one throughout.
Figure~\ref{fig:mc-crossover} plots the widths and the projection distance.

\tabMcCrossover

\figMcCrossover

The MDS route loses calibration through its studentization. It reconstructs a level
covariance from the pairwise standard errors and projects it to positive
semi-definiteness, but still studentizes the simulated contrasts by the original
$\seNW{jk}$. When the standard-error matrix is badly non-Euclidean, the projection
distorts the implied pairwise variances. The studentized draws are then no longer
unit-variance, the maximum in \eqref{eq:cv-sim} inflates, and the critical value, hence
every confidence set, balloons. The direct construction studentizes by
$\sqrt{\Omegahat^{+}_{aa}}$, the diagonal of its own projected matrix, so its draws are
unit-variance by construction.

The real data confirm the mechanism at its extreme. On the unbalanced, heavy-tailed NGDP
panel (Chapter~\ref{ch:apps}), the MDS studentized draws have standard deviations ranging
from $1.0$ to $12.1$ and a two-sided $5\%$ critical value of $24.1$, which collapses
every rank set to $[1,p]$. The direct construction's draws are unit-variance, and its
$5\%$ critical value is $2.7$. The Monte Carlo evidence thus has two parts. Under the
idealized model the two constructions are validated and asymptotically equivalent. The
direct construction's advantage is a robustness property that appears precisely when
non-stationarity and imbalance drive the standard-error matrix away from Euclidean
embeddability.
